\subsection{Bounded ordinary-Hecke complexes and the nonstandard quadratic relation}

The Euler lift in \fid{F27} does not supply a categorical quadratic
relation. The following test allows arbitrary bounded differentials in
the complementary term and uses the actual required relation.

Work over a field \(k\) of characteristic zero with an ordinary faithful
reflection realization of a finite Coxeter group \(W\). The standard
type-A realization over \(\mathbb Q\) suffices in every rank at issue.
Let \(R\) be the polynomial ring on the realization, with linear
coordinates in internal degree two, and let \(\mathcal B\) be its
ordinary Soergel category. The conventions are
\[
 B_s=R\otimes_{R^s}R(1),\qquad
 M(1)^d=M^{d+1},\qquad b_s=[B_s],\quad
 \delta=v+v^{-1},\quad c_s=b_s-\delta.
\]
Internal shift \((1)\) and cohomological shift \([1]\) are distinct.
These are the conventions of \cite[\S\S3.3--3.5,
pp.~1105--1108]{EW}. The identities
\(b_s^2=\delta b_s\), \(c_s^2=-\delta c_s\), and
\(b_sc_s=0\) are initially identities of classes.

Use two independent copies of the realization and put
\(A=R\otimes_kR\). In the bounded homotopy category of
Soergel bimodules for \(W\times W\), set
\[
 s_1=(s,e),\qquad s_2=(e,s),\qquad u=(s,s),
\]
\[
 P_s=b_s\otimes b_s+c_s\otimes c_s,
 \qquad
 \mathscr D=\mathbf1(2)\oplus\mathbf1^{\oplus2}
                        \oplus\mathbf1(-2).
\]
Thus \([\mathscr D]=\delta^2\), and \(P_s\) is the actual
named nonstandard element, with its original ambient inclusion.

\begin{finding}[\fid{F59}: a direct-sum derived lift cannot satisfy the quadratic]
\label{f:soergelderivedquadratic}
There is no bounded complex \(X_s\) with Euler class \(P_s\),
an isomorphism
\begin{equation}\label{eq:soergelactualquadratic}
 X_s\otimes_A X_s\simeq\mathscr D\otimes_A X_s,
\end{equation}
and a direct-sum decomposition
\begin{equation}\label{eq:soergelsplithypothesis}
 X_s\simeq(B_s\boxtimes B_s)\oplus Y_s,
 \qquad [Y_s]=c_s\otimes c_s.
\end{equation}
In particular, any construction
\[
 (B_s\boxtimes B_s)\oplus
 (C_s^{(1)}\boxtimes C_s^{(2)})
\]
with bounded ordinary-Hecke complexes of Euler class
\([C_s^{(j)}]=c_s\) fails
\eqref{eq:soergelactualquadratic}, regardless of the two cone
choices and their differentials.

For any bounded lift satisfying the quadratic, without the
direct-sum assumption, ungraded generic localization must have
cohomology only in degree zero, equal to
\[
                    Q_e^{\oplus2}\oplus Q_u^{\oplus2},
 \qquad Q=\operatorname{Frac}(A).
\]
This is a necessary condition. It does not exclude every
coupled bounded complex.
\end{finding}

\begin{proof}
\par\smallskip\noindent\textbf{Generic graph bimodules.}\quad
After localization and forgetting the internal grading, the
ordinary category maps monoidally to finite sums of graph
bimodules \(Q_w\), \(w\in W\times W\). Here \(Q_w\)
is \(Q\) as a left vector space and has right action
\(z\cdot a=w(a)z\). The corrected primary localization
framework is \cite[\S1.1, pp.~3--4; \S1.2,
Theorem~1.6 and Remark~1.7, pp.~4--5;
\S2.1, equations~(2.1)--(2.6), p.~7]{EWLocalized}.
The elementary facts needed here can also be checked directly.
Faithfulness gives
\begin{equation}\label{eq:soergelgraphfacts}
 \operatorname{Hom}_{Q-Q}(Q_w,Q_z)=
 \begin{cases}Q,&w=z,\\0,&w\ne z,\end{cases}
 \qquad Q_w\otimes_Q Q_z\simeq Q_{wz}.
\end{equation}
The second map sends \(a\otimes b\) to \(a\,w(b)\).
A nonzero map in the first formula for distinct labels would
be a scalar killed by a nonzero \(w(a)-z(a)\) in a field.
Thus the additive graph category is semisimple, and the
tensor product of two nonzero objects is nonzero.

For a simple reflection the splitting is
\[
 Q\otimes_{Q^s}Q\longrightarrow Q_e\oplus Q_s,
 \qquad a\otimes b\longmapsto(ab,a\,s(b)).
\]
If \(s(\alpha)=-\alpha\), its two standard vectors are
the images of
\[
       \frac{1\otimes1\pm\alpha^{-1}\otimes\alpha}{2}.
\]
Tensor products give the decomposition for Bott--Samelson
bimodules, and direct summands remain sums of graph lines.
Denominators move through a graph line by its automorphism,
so this is a monoidal localization. Denote its extension
to bounded complexes by \(\Lambda\).

\par\smallskip\noindent\textbf{A generic amplitude lemma.}\quad
Let \(Z\) be a nonzero bounded graph complex satisfying
\(Z\otimes_Q Z\simeq Z^{\oplus m}\) with \(m>0\).
Semisimplicity splits each bounded complex into its
cohomology and contractible pairs. Tensoring is exact and
preserves contractibility. If \(p,r\) are its least and
greatest nonzero cohomology degrees, the square has
extremes \(2p,2r\): the extreme tensor powers of
\(H^pZ,H^rZ\) are nonzero by
\eqref{eq:soergelgraphfacts}. The right side has extremes
\(p,r\), hence \(2p=p\) and \(2r=r\). Thus \(p=r=0\).

The graph Euler ring is \(\mathbb Z[W\times W]\), and
\[
 [\Lambda B_s]=e+s,\qquad
 [\Lambda c_s]=s-e,\qquad
 [\Lambda P_s]=2(e+u).
\]
The middle notation denotes the image of the virtual
Hecke class. Internal shifts have been forgotten, so the
coefficient \(\mathscr D\) becomes four unit copies.
The localized \(X_s\) is nonzero because its Euler class
is \(2(e+u)\). Its quadratic therefore forces it into
degree zero by the amplitude lemma.

Under \eqref{eq:soergelsplithypothesis}, the direct
summand \(\Lambda Y_s\) also has cohomology only in degree
zero. Its Euler graph multiplicities must be nonnegative.
But its required class is
\[
             (s_1-e)(s_2-e)=e+u-s_1-s_2.
\]
The four labels are distinct, so the negative coefficients
give a contradiction. This proves the direct-sum assertion
independently of every differential chosen inside \(Y_s\).
Without that assertion, the same lemma and the Euler class
give precisely the stated necessary localized cohomology.
\end{proof}

\par\smallskip\noindent\textbf{An explicit unit cone and its failed square.}\quad
The natural maps in one factor are
\[
 \eta:R(-1)\longrightarrow B_s,
 \quad f\longmapsto f(\alpha\otimes1+1\otimes\alpha),
 \qquad
 m:B_s\longrightarrow R(1),\quad f\otimes g\longmapsto fg.
\]
They preserve the displayed grading. The decomposition
\(R=R^s\oplus\alpha R^s\) gives the exact sequences
\begin{equation}\label{eq:soergelcounitexact}
 0\longrightarrow R_s(-1)
 \xrightarrow{\ f\mapsto f(1\otimes\alpha-\alpha\otimes1)\ }
 B_s\xrightarrow{m}R(1)\longrightarrow0,
\end{equation}
\begin{equation}\label{eq:soergelunitexact}
 0\longrightarrow R(-1)\xrightarrow{\eta}B_s
 \xrightarrow{\ f\otimes g\mapsto f\,s(g)\ }R_s(1)
 \longrightarrow0.
\end{equation}
To check them, use the left basis
\(1\otimes1,1\otimes\alpha\). The two target maps are
surjective; their rank-one kernels are respectively the
anti-invariant and invariant displayed vectors. The domain
property gives injectivity of the first arrows. These are
the standard dot and multiplication maps of
\cite[\S3.4, pp.~1106--1107]{EW}; the exact sequences have
just been verified explicitly with our shifts.

Choose the complementary cone
\[
 C_s=\operatorname{Cone}\bigl(
 R(-1)\oplus R(1)\xrightarrow{(\eta,0)}B_s\bigr),
\]
with the source in degree minus one and target in degree
zero. Its Euler class is \(c_s\). Every degree-zero map
from these two unit shifts to \(B_s\) has this form up to
a scalar: a central element of \(B_s\) is a polynomial
multiple of \(\alpha\otimes1+1\otimes\alpha\), whose
minimum degree is one. Its degree-one space is one
dimensional and there is no degree-minus-one element.
The zero scalar is already covered by the finding.

In the derived category of bimodules,
\eqref{eq:soergelunitexact} gives
\[
                     C_s\simeq R_s(1)\oplus R(1)[1].
\]
Thus the direct-sum recipe is isomorphic there to
\begin{equation}\label{eq:soergeldirectconeimage}
 (B_s\boxtimes B_s)\oplus A_u(2)
 \oplus A_{s_1}(2)[1]\oplus A_{s_2}(2)[1]
 \oplus A(2)[2].
\end{equation}
All modules in this expression are free on either side.
Its minimum nonzero cohomological degree is minus two,
with cohomology \(A(2)\); its square has nonzero
degree-minus-four cohomology \(A(4)\).
Tensoring \eqref{eq:soergeldirectconeimage} with
\(\mathscr D\) produces no degree-minus-four cohomology.
Consequently the natural specified cone constructs an
Euler lift but fails the quadratic even before generic
localization. The finding covers other cone conventions
without assuming their differentials are zero.

\par\smallskip\noindent\textbf{A coupled Rouquier Euler lift.}\quad
There is a different concrete lift that passes the generic
amplitude test. Put
\[
 F_s=[B_s\xrightarrow{m}R(1)],\qquad
 E_s=[R(-1)\xrightarrow{\eta}B_s],
\]
where \(B_s\) is in degree zero in both complexes. These
are the positive and inverse Rouquier complexes in this
normalization; see \cite[\S6.4, p.~1121]{EW} and
\cite[\S3.2.4 and Lemma~3.3, p.~9; \S3.3.1,
p.~10]{RouquierBraid}, with its different raw grading
normalization translated by our displayed maps.
Define
\begin{equation}\label{eq:soergelcoupledRouquier}
 X_s^{\rm Rouq}=(F_s\boxtimes F_s)
       \oplus(E_s\boxtimes E_s)\oplus A^{\oplus2}.
\end{equation}
Its exact Euler class is
\begin{align*}
 &(b_s-v)\otimes(b_s-v)
 +(b_s-v^{-1})\otimes(b_s-v^{-1})+2\\
 &\qquad=2b_s\otimes b_s
    -\delta(b_s\otimes1+1\otimes b_s)+\delta^2=P_s.
\end{align*}
The two scalar constant terms in the tensor products are
\(v^2,v^{-2}\), and the two unit summands supply the
remaining constant two. Thus this preserves the original
ambient class rather than changing its inclusion as in
\fid{F16}.

The exact sequences give
\[
          F_s\simeq R_s(-1),\qquad E_s\simeq R_s(1)
\]
in the derived category of bimodules. Consequently
\begin{equation}\label{eq:soergelRouquierderivedimage}
 X_s^{\rm Rouq}\simeq A^{\oplus2}
                       \oplus A_u(-2)\oplus A_u(2).
\end{equation}
This has cohomology only in degree zero, and after
ungraded generic localization has exactly the support
required in the finding. A generic-amplitude argument
alone cannot exclude all coupled Euler lifts.

Its actual graded quadratic nevertheless fails. The graph
modules \(A\) and \(A_u\) have no nonzero bimodule map
between them: such a map is a scalar killed by a nonzero
\(a-u(a)\) in the polynomial domain. An isomorphism
between sums of their shifts must preserve the two graph
summands and their graded free ranks.
The identity-graph coefficient in the square of
\eqref{eq:soergelRouquierderivedimage} is
\[
               4+(v^2+v^{-2})^2=v^4+6+v^{-4},
\]
whereas its coefficient in
\(\mathscr D\otimes_A X_s^{\rm Rouq}\) is
\[
                         2\delta^2=2v^2+4+2v^{-2}.
\]
These are unequal. All graph modules are free on either
side, so no higher Tor terms have been omitted in this
comparison. The two sides are not isomorphic even in
the derived category of graded bimodules, and therefore
not in the bounded ordinary-Hecke homotopy category.
This calculation uses the specified Rouquier maps and
exact sequences, not an assumption that ordinary braid
coherence implies the nonstandard quadratic.

The direct-sum failure is a necessary one-generator test
in every rank. It is stronger than \fid{F27}'s
zero-differential amplitude calculation, but its
direct-sum premise remains explicit. At the generic
ungraded level the object
\(Q_e^{\oplus2}\oplus Q_u^{\oplus2}\) does satisfy the
fourfold quadratic, by the multiplication table of the
two-element subgroup \(\{e,u\}\).
The remaining problem is descent of a coupled finite
graded complex with the actual quadratic isomorphism,
followed by coherent higher nonstandard relation maps.
No result here supplies VIII's designated lattices or
an arithmetic Frobenius comparison.

% Separate paragraph for Part III.
\paragraph{\fid{F59}'s bounded-Hecke test in Grothendieck's language.}
The ordinary bimodules, finite reflection correspondences
and graph restrictions are quasi-coherent affine module
data, through EGA~I, Theorem~1.3.7, p.~86,
Corollary~1.3.8, p.~87, and Theorem~1.4.1,
pp.~90--92 \cite{EGA}.
The bounded complexes have terms free on each tensor
side. The perfect-complex vocabulary is SGA~6,
I, Definition~4.7, pp.~121--122, and Example~4.8,
p.~122, with that tensor base retained explicitly
\cite{SGA6}.
The graph tensor category remembers group order;
no commutative \(\lambda\)-ring is asserted for its
noncommutative Grothendieck multiplication.
SGA~6, V, \S\S2--3 and cote~133, scan~18 provide
the coefficient and virtual-class language, while
SGA~3, I, 4.7.2--3 and cote~90, scan~17 provide
the action framework \cite{SGA3,SGA6,Archive}.
No group-scheme realization of the noncommutative algebra
follows from that vocabulary.
The new proof separates an Euler identity from an actual
relation isomorphism; its effective-multiplicity
contradiction uses precisely the direct-sum premise.
Generic restriction is a necessary test rather than a
descent theorem. The new cone, localization and grading
calculations are proved above and are not attributed to
the archive manuscripts.
